\section{Experimental Setup}
\label{sec:setup}

All experiments are run with the same decoder-only Transformer training code based on the nanoGPT implementation(\cite{karpathy_nanogpt}) and configured through Hydra.
Our goal is not to match the largest models studied in the original paper, but
to reproduce the same failure modes in a compact setting that is cheap to run
and easy to ablate.

\paragraph{Model.}
We closely follow the paper's implementation : The model uses rotary positional embeddings, no dropout, bias-free linear layers, and no weight tying between the token embedding matrix and the output
projection.

\paragraph{Data and optimisation.}
Due to size constraints we chose to train on the OpenWebText-style dataset provided in our pipeline, using the
standard next-token prediction objective and not on the C4 dataset.
The default batch size is 32 sequences of length 256, with one gradient
accumulation step.
Training is performed in bfloat16 when available, with \texttt{torch.compile}
enabled.
We use AdamW with $\beta_1=0.9$ and $\beta_2=0.95$.
